\documentclass[aps,prd,twocolumn,superscriptaddress,nofootinbib,showkeys]{revtex4-2}
\usepackage{amsmath,amssymb,hyperref}
\begin{document}
\title{CERN+1 (formerly Janus-LHC): Wino Pair Cross Sections at Any LHC Energy from Drell--Yan Scaling with One Common Evolution Exponent}
\author{Danny Eldridge}
\thanks{CERN+1 is independent work. It is not affiliated with, sponsored by or endorsed by CERN or any LHC experiment; the name means ``the chain, plus one more check''. Version 1 of this note used the working name Janus-LHC.}
\email{dannyeldridge54@gmail.com}
\affiliation{Independent researcher, Oklahoma City, Oklahoma, USA}
\date{18 November 2026 (v2, referee corrections)}
\begin{abstract}
Searches for electroweakinos at the LHC need production cross sections at the collider's actual energy. Reference values are computed with Resummino at aNNLO+NNLL accuracy, one mass point at a time, and are tabulated at 13 and 13.6~TeV~\cite{Res}. We present the cross-section module of CERN+1 (formerly Janus-LHC), a small open toolkit with two routes to these numbers. The first interpolates the tables in $\ln\sigma$ against $\ln m$. A leave-one-out test gives a median error of 0.015--0.016\% and a maximum of 0.30\%. The second uses no parton-distribution (PDF) file directly; the PDFs enter only through the 13~TeV reference table. At leading order, if the parton densities did not evolve with scale and the pair mass were the only scale, $s\,\sigma$ would depend only on $\tau_0=4m^2/s$, the classic Drell--Yan scaling~\cite{DY}. With the scaling function $\Phi(\tau_0)$ built from the 13~TeV table alone, this route predicts every 13.6~TeV entry for wino-like $\tilde\chi_2^0\tilde\chi_1^\pm$ and $\tilde\chi_1^+\tilde\chi_1^-$ (200--2000~GeV) to $\le1.0\%$. One correction factor $(\sqrt s/13\,\mathrm{TeV})^{\kappa}$, fitted to the same 13.6~TeV entries, brings every entry to $\le0.29\%$. The fitted exponent is nearly the same for the three channels, $\kappa=-0.161$, $-0.163$, $-0.165$, as expected for $q\bar q'$ channels that share the same parton luminosities; it absorbs mainly the scale evolution of the parton densities. At 14~TeV (HL-LHC) the two routes agree to 0.7--1.0\% from 500 to 1800~GeV, so we quote their mean with the spread as the error. This error is interpolation and route spread only; it does not include the PDF and scale uncertainties of the reference tables, which are larger and must be added from Ref.~\cite{Res}. The toolkit also sets 95\% CLs limits for counting experiments in two ways, with pyhf (asymptotic) and with exact Poisson sums; they differ by up to 8\% at low counts, and the toolkit uses the exact result there. The three cross-section tests were pre-registered.
\end{abstract}
\keywords{LHC; electroweakino; wino; cross section; Drell--Yan scaling; parton luminosity; HL-LHC; CLs; pyhf; reinterpretation; reproducibility}
\maketitle

\section{Motivation}
Supersymmetric and minimal dark-matter models with a wino-like state need the electroweakino pair cross section at whatever energy a search uses: 13, 13.6 or 14~TeV. The standard values come from Resummino at NLO+NLL or aNNLO+NNLL, with global parton distributions and a separate run for each mass and energy~\cite{Res}. We wanted three things: an answer at any mass and energy in microseconds, a stated numerical error on that answer, and a check that does not call a parton-distribution file. That file could not be downloaded in our environment. Rescaling cross sections between collider energies with parton-luminosity ratios is standard practice~\cite{Quigg}; our contribution is a calibrated one-parameter version of it on the Resummino wino tables, with a pre-registered test.

\section{Route 1: interpolation of the reference tables}
We take the aNNLO+NNLL tables for mass-degenerate winos of Ref.~\cite{Res} (Table~6), for $\tilde\chi_2^0\tilde\chi_1^+$, $\tilde\chi_2^0\tilde\chi_1^-$ and $\tilde\chi_1^+\tilde\chi_1^-$, from 100 to 2000~GeV in steps of 100~GeV, at 13 and 13.6~TeV. The interpolation is a cubic spline in $\ln\sigma$ against $\ln m$, because $\sigma$ is close to a power law in $m$. We tested it by dropping each interior point and predicting it from the rest (18 masses $\times$ 3 channels $=54$ points per energy). The median error is 0.015\% at 13~TeV and 0.016\% at 13.6~TeV. The maximum is 0.30\% at 13~TeV and 0.15\% at 13.6~TeV. A spline in $\ln\sigma$ against $m$ gives a 10.3\% maximum error near 200~GeV, so the $\ln m$ axis matters. This error measures interpolation only; the tables' own scale and PDF uncertainties are separate and larger.

\section{Route 2: Drell--Yan scaling}
At leading order, $\sigma=\sum_{ij}\int d\tau\,\mathcal L_{ij}(\tau;\mu)\,\hat\sigma_{ij}(\tau s)$. If the partonic cross section depends on the pair mass $m$ as its only scale, $\hat\sigma_{ij}(\hat s)=f_{ij}(4m^2/\hat s)/\hat s$ (this neglects $M_{W,Z}^2/\hat s$ terms in the propagators and any squark exchange), and if the parton luminosity did not depend on the factorisation scale $\mu$, then $s\,\sigma$ would be a function of $\tau_0=4m^2/s$ alone~\cite{DY}. We parametrise the departure from this as
\begin{equation}
s\,\sigma(m,\sqrt s)=\Phi(\tau_0)\,\Bigl(\frac{\sqrt s}{13\ \mathrm{TeV}}\Bigr)^{\kappa}.
\label{eq:scal}
\end{equation}
Here $\kappa=0$ is exact scaling. Because $\mu\sim m=\tfrac12\sqrt{\tau_0 s}$ grows with $\sqrt s$ at fixed $\tau_0$, $\kappa\neq0$ absorbs mainly the scale evolution of the parton densities; it also absorbs the running of $\alpha_s$ in the higher-order corrections and the residual electroweak-mass effects. We built $\Phi$ from the 13~TeV table only (a cubic spline in $\ln\tau_0$) and predicted the 13.6~TeV table. Before computing, we pre-registered three predictions~\cite{P4}:
\begin{itemize}
\item[(P1)] exact scaling is good to 5\%;
\item[(P2)] one exponent per process gives $\le1\%$;
\item[(P3)] the two routes agree at 14~TeV to 5\%.
\end{itemize}

\emph{Results.} All three passed. P1 and P2 use the 19 masses from 200 to 2000~GeV per channel.
\begin{itemize}
\item \emph{P1.} With $\kappa=0$, the deviations lie between $+0.45\%$ and $+1.0\%$ for all three processes. They have one sign and vary smoothly with mass.
\item \emph{P2.} Fitting $\kappa$ by least squares to the same 13.6~TeV entries gives $-0.165$, $-0.163$ and $-0.161$ for $\tilde\chi_2^0\tilde\chi_1^+$, $\tilde\chi_2^0\tilde\chi_1^-$ and $\tilde\chi_1^+\tilde\chi_1^-$. The maximum deviation then falls to 0.29\%, 0.26\% and 0.26\%. This is an in-sample, one-parameter fit, not an out-of-sample prediction.
\item \emph{P3.} At 14~TeV, Eq.~(\ref{eq:scal}) and a log-linear-in-$\sqrt s$ extrapolation of the two tables agree to 0.71--0.96\% at the five masses tested (500, 1000, 1100, 1500 and 1800~GeV). The scaling route is the lower one at every point.
\end{itemize}
That $\kappa$ is nearly the same for the three processes is consistent with its origin in parton evolution: these $q\bar q'$ channels are produced from closely related quark--antiquark luminosities. With only two energies, a lever arm of 4.6\% in $\sqrt s$, $\kappa$ is an effective slope, and its constancy beyond 13.6~TeV is an assumption.

\emph{Example.} For $m=1.1$~TeV at 14~TeV, $\sigma(\tilde\chi_2^0\tilde\chi_1^+)=0.787$~fb $\pm0.8\%$. The $\pm0.8\%$ is the full spread of the two routes. Both routes are two-point extrapolations from the same tables, so the spread is a measure of extrapolation ambiguity, not a statistically independent error. It does not include the PDF and scale uncertainties quoted with the reference tables~\cite{Res}, which must be added to it.

\section{Statistics}
For counting experiments the toolkit computes 95\% CLs upper limits~\cite{Junk,Read} on the signal yield in two ways: with pyhf~\cite{pyhf}, using the asymptotic distribution of $\tilde q_\mu$~\cite{CCGV} and a negligible background uncertainty, and with exact Poisson sums, $\mathrm{CLs}=P(N\le n\,|\,s+b)/P(N\le n\,|\,b)$, which use the observed count as the test statistic and no background uncertainty. Pairs (pyhf vs exact): for $b=0.5$, $n=0$, 3.07 vs 3.00; for $b=3$, $n=3$, 5.00 vs 5.40; for $b=10$, $n=12$, 9.51 vs 9.97. The two differ by up to 8\% at low counts (relative to the pyhf value). The difference reflects both the asymptotic approximation and the different test statistic. The toolkit reports the exact value at low counts. These comparisons were not part of the pre-registered tests.

\section{Limitations and next steps}
Both routes inherit the central values and the precision of the reference tables, and neither propagates their PDF or scale uncertainties; users must add those from Ref.~\cite{Res}. Route~2 is calibrated on 13 and 13.6~TeV and extrapolates to 14~TeV, which is a small step in $\tau_0$; far larger extrapolations in $\sqrt s$ are not covered by this test. Next: an independent leading-order computation with a global PDF set, giving a K-factor against Resummino; a disappearing-track recast from published HEPData efficiency maps; and published full likelihoods in pyhf. Code: \texttt{tools/lhc} in the archive~\cite{AEGIS}.

\emph{Method note.} Prepared with an AI assistant (Claude, Anthropic) under the workflow of Ref.~\cite{P4}. Note D1299 and lock \texttt{PREREG\_LHC1\_scaling} in the archive. Version 2 renames the toolkit, states that the quoted errors exclude PDF and scale uncertainties, and adds the prior-art and statistics references.

\emph{License.} \copyright~2026 Danny Eldridge. Free for individuals, scientists and students under CC BY-NC-SA 4.0 (AEGIS Community License); any business or for-profit use requires the AEGIS Commercial License (see LICENSE in the archive).

\begin{thebibliography}{99}
\bibitem{Res} J.~Fiaschi, B.~Fuks, M.~Klasen and A.~Neuwirth, ``Electroweak superpartner production at 13.6 TeV with Resummino,'' Eur.\ Phys.\ J.\ C \textbf{83}, 707 (2023), arXiv:2304.11915.
\bibitem{DY} S.~D.~Drell and T.-M.~Yan, ``Massive lepton-pair production in hadron-hadron collisions at high energies,'' Phys.\ Rev.\ Lett.\ \textbf{25}, 316 (1970).
\bibitem{Quigg} C.~Quigg, ``LHC physics potential vs.\ energy,'' arXiv:0908.3660 (2009).
\bibitem{Junk} T.~Junk, ``Confidence level computation for combining searches with small statistics,'' Nucl.\ Instrum.\ Meth.\ A \textbf{434}, 435 (1999).
\bibitem{Read} A.~L.~Read, ``Presentation of search results: the CL$_s$ technique,'' J.\ Phys.\ G \textbf{28}, 2693 (2002).
\bibitem{CCGV} G.~Cowan, K.~Cranmer, E.~Gross and O.~Vitells, ``Asymptotic formulae for likelihood-based tests of new physics,'' Eur.\ Phys.\ J.\ C \textbf{71}, 1554 (2011).
\bibitem{pyhf} L.~Heinrich, M.~Feickert, G.~Stark and K.~Cranmer, ``pyhf: pure-Python implementation of HistFactory statistical models,'' J.\ Open Source Softw.\ \textbf{6}, 2823 (2021).
\bibitem{P4} D.~Eldridge, ``Lock Before You Compute: A Pre-Registration Workflow for Theoretical Physics Done with an AI Assistant,'' Zenodo (2026), doi:10.5281/zenodo.24291026.
\bibitem{AEGIS} Reproducibility archive: \url{https://github.com/dannyeldridge54/aegis}.
\end{thebibliography}
\end{document}
